\section{Podstawowe pojęcia}
Zanim przystąpimy do konstrukcji narzędzi mierzących skomplikowane zbiory, musimy precyzyjnie określić środowisko, w którym będziemy operować. Nie możemy oczekiwać, że każda dowolna rodzina podzbiorów zachowa się stabilnie przy operacjach sumowania czy dopełniania. W niniejszym rozdziale wprowadzamy zatem algebraiczny szkielet teorii — od pierścieni i ciał, przez kluczowe pojęcie $\sigma$-algebry, aż po topologię euklidesową przestrzeni $\Rn$, która wyznaczy geometryczny punkt odniesienia dla pojęcia objętości.

\subsection{Struktury rodzin zbiorów}
W poniższych definicjach niech $X$ będzie ustalonym niepustym zbiorem.
\begin{defi}[Pierścień zbiorów]
    \textbf{Pierścieniem podzbiorów} przestrzeni $X$ nazywamy rodzinę $\mathcal{R} \subs \mathcal{P}(X)$, zawierającą $\emptyset$, zamkniętą na skończone sumy i różnice zbiorów, tj. spełniającą warunki:
    \begin{enumerate}
        \item $\emptyset \in \mathcal{R}$
        \item $\forall_{A, B \in \mathcal{R}} A \cup B \in \mathcal{R}$
        \item $\forall_{A, B \in \mathcal{R}} A \setminus B \in \mathcal{R}$
    \end{enumerate}
\end{defi}

\begin{defi}[Algebra (ciało) zbiorów]
    \textbf{Algebrą (bądź ciałem)} podzbiorów przestrzeni $X$  nazywamy pierścień $\mathcal{F}$ podzbiorów $X$, zawierający całą przestrzeń $X$. Równoważnie, algebrę podzbiorów przestrzeni $X$ możemy zdefiniować jako rodzinę $\mathcal{F}$ podzbiorów przestrzeni $X$ spełniającą następujące warunki:
    \begin{enumerate}
        \item $\emptyset \in \mathcal{F}$
        \item $\forall_{A,B \in \mathcal{F}} A \cup B \in \mathcal{F}$
        \item $\forall_{A \in \mathcal{F}} X \setminus A \in \mathcal{F}$
    \end{enumerate}
\end{defi}

\begin{defi}[$\sigma$-algebra ($\sigma$-ciało)]
    \textbf{$\sigma$-ciałem (lub przeliczalnie addytywnym ciałem)} podzbiorów przestrzeni $X$ nazywamy rodzinę $\mathcal{M}$ podzbiorów $X$ spełniającą następujące warunki:
    \begin{enumerate}
        \item $\emptyset \in \mathcal{M}$
        \item $\forall_{A \in \mathcal{M}} X \setminus A \in \mathcal{M}$
        \item $\forall_{A_1, A_2, \ldots \in \mathcal{M}} \bigcup_{n = 1}^\infty A_n \in \mathcal{M}$
    \end{enumerate}
    Parę $(X, \mathcal{M})$ nazywamy \textbf{przestrzenią mierzalną}, a zbiory należące do rodziny $\mathcal{M}$ zbiorami $\mathcal{M}$-mierzalnymi (lub mierzalnymi).
    $\sigma$-algebra podzbiorów przestrzeni $X$ jest więc algebrą w której ,,warunek sumy'' rozszerzono na sumę przeliczalnej rodziny zbiorów. W oczywisty sposób każda $\sigma$-algebra jest algebrą, choć odwrotne twierdzenie nie zawsze jest prawdziwe.
\end{defi}

\begin{defi}[$\sigma$-algebra generowana]
    $\sigma$-algebrą generowaną przez rodzinę zbiorów $\mathcal{F}$ nazywamy najmniejszą $\sigma$-algebrę zawierającą tę rodzinę i oznaczamy $\sigma(\mathcal{F})$.
\end{defi}

\subsection{Elementarna geometria i topologia w $\Rn$}

W całej pracy przestrzeń $\Rn$ ($n \in \N$, $n \wlr 1$) traktujemy jako $n$-wymiarową przestrzeń liniową nad ciałem liczb rzeczywistych $\R$. Punkty przestrzeni oznaczamy przez $x = (x_1, \ldots, x_n)$, $y = (y_1, \ldots, y_n) \in \Rn$.

\begin{defi}[Metryka i topologia euklidesowa]
    Dla dowolnych punktów $x, y \in \Rn$ \textbf{odległością euklidesową} (metryką euklidesową) nazywamy liczbę
    \begin{equation*}
        d(x, y) = \|x - y\| = \sqrt{\sum_{i=1}^n (x_i - y_i)^2}.
    \end{equation*}
    \begin{enumerate}
        \item \textbf{Kulą otwartą} o środku w punkcie $x_0 \in \Rn$ i promieniu $r > 0$ nazywamy zbiór
              \begin{equation*}
                  B(x_0, r) := \setof{x \in \Rn}{d(x, x_0) < r}
              \end{equation*}

        \item Zbiór $U \subseteq  \Rn$ nazywamy \textbf{zbiorem otwartym}, jeżeli dla każdego $x \in U$ istnieje $r > 0$ takie, że $B(x, r) \subseteq  U$.
              Rodzinę wszystkich zbiorów otwartych w $\Rn$ oznaczamy przez $\mathrm{top}(\Rn)$.

        \item Zbiór $F \subseteq  \Rn$ nazywamy \textbf{zbiorem domkniętym}, jeżeli jego dopełnienie $F^c = \Rn \setminus F$ jest zbiorem otwartym ($F^c \in \mathrm{top}(\Rn)$).
              Rodzinę wszystkich zbiorów domkniętych oznaczamy przez $\mathrm{cotop}(\Rn)$.

        \item Zbiór $K \subseteq  \Rn$ nazywamy \textbf{zbiorem zwartym}, jeżeli z każdego jego pokrycia zbiorami otwartymi można wybrać podpokrycie skończone. Na mocy twierdzenia Heinego-Borela zbiór $K \subseteq  \Rn$ jest zwarty wtedy i tylko wtedy, gdy jest domknięty i ograniczony w metryce euklidesowej.

        \item Dla dowolnego niepustego zbioru $A \subseteq \Rn$ liczbę
              \begin{equation*}
                  \diam(A) := \sup\setof{d(x, y)}{x, y \in A} \in [0, +\infty]
              \end{equation*} nazywamy \textbf{średnicą zbioru} $A$. Przyjmujemy ponadto $\diam(\emptyset) = 0$.
        \item Dla dwóch niepustych podzbiorów $A, B \subseteq \Rn$ \textbf{odległością zbiorów} $A$ i $B$ nazywamy wielkość
              \begin{equation*}
                  \dist(A, B) := \inf\setof{d(x, y)}{x \in A,  y \in B}.
              \end{equation*}
              Przyjmujemy $\dist(A, B) = +\infty$, gdy $A = \emptyset$ lub $B = \emptyset$. Dla punktu $x \in \Rn$ piszemy $\dist(x, A) := \dist(\{x\}, A)$.
    \end{enumerate}
\end{defi}

\begin{defi}[Kostka w $\Rn$]
    Niech $a_j, b_j \in \R$ spełniają $a_j < b_j$ dla każdego $j \in \{1, \ldots, n\}$.

    \begin{enumerate}
        \item \textbf{Kostką domkniętą} (przedziałem wielowymiarowym domkniętym) w $\Rn$ nazywamy iloczyn kartezjański odcinków domkniętych:
              \begin{equation*} P = [a_1, b_1] \times [a_2, b_2] \times \ldots \times [a_n, b_n] = \prod_{j=1}^n [a_j, b_j].
              \end{equation*}
        \item \textbf{Kostką otwartą} w $\Rn$ nazywamy zbiór postaci:
              \begin{equation*} Q = (a_1, b_1) \times (a_2, b_2) \times \ldots \times (a_n, b_n) = \prod_{j=1}^n (a_j, b_j).
              \end{equation*}
        \item Ogólnie, \textbf{kostką ograniczoną} (przedziałem $n$-wymiarowym) w $\Rn$ nazywamy zbiór będący iloczynem kartezjańskim dowolnych przedziałów jednowymiarowych o końcach $a_j, b_j$ (domkniętych, otwartych lub jednostronnie domkniętych). Każda taka kostka $R$ spełnia relację $Q \subseteq  R \subseteq  P$ dla odpowiednich kostek otwartej $Q$ i domkniętej $P$.
    \end{enumerate}
\end{defi}

\begin{defi}[Elementarna objętość kostki] Dla dowolnej ograniczonej kostki $R \subseteq \Rn$ (domkniętej, otwartej bądź częściowo domkniętej) o końcach odcinków krawędziowych $a_j < b_j$ ($j = 1, \ldots, n$) definiujemy jej \textbf{objętość elementarną} wzorem:
    \begin{equation*}
        V^n(R) := \prod_{j=1}^n (b_j - a_j).
    \end{equation*}
    Ponadto dla zbioru pustego kładziemy $V^n(\emptyset) := 0$.
\end{defi}

\begin{defi}[$\sigma$-algebra zbiorów borelowskich]
    Niech $\mathrm{top}(\Rn)$ oznacza rodzinę wszystkich podzbiorów otwartych w $\Rn$. $\sigma$-algebrę generowaną przez tę rodzinę, oznaczaną symbolem $\BorelRn$, nazywamy \textbf{$\sigma$-algebrą zbiorów borelowskich} w $\Rn$:
    \begin{equation*}
        \Borel{\Rn} := \sigma(\mathrm{top}(\Rn)).
    \end{equation*}
    Elementy rodziny $\BorelRn$ nazywamy \textbf{zbiorami borelowskimi}.
\end{defi}
\begin{uwaga}
    Zauważmy, że ponieważ dopełnienie zbioru otwartego jest zbiorem domkniętym, zachodzi tożsamość: \begin{equation*} \BorelRn = \sigma(\mathrm{cotop}(\Rn)).
    \end{equation*}
    Co więcej, $\sigma$-algebra $\BorelRn$ jest generowana przez każdą z poniższych rodzin:
    \begin{enumerate}
        \item rodzinę wszystkich kostek otwartych w $\Rn$;
        \item rodzinę wszystkich kostek domkniętych w $\Rn$;
        \item rodzinę kostek postaci $\prod_{j=1}^n [a_j, b_j)$ lub $\prod_{j=1}^n (a_j, b_j)$ o końcach wymiernych $a_j, b_j \in \Q$.
    \end{enumerate}
\end{uwaga}
\begin{defi}[Klasy borelowskie $G_\delta$ oraz $F_\sigma$]
    Niech $Y \subseteq  \Rn$.
    \begin{enumerate}
        \item Zbiór $Y$ nazywamy \textbf{zbiorem typu $G_\delta$}, jeżeli jest on przeliczalnym przekrojem zbiorów otwartych, to znaczy:
              \begin{equation*}
                  \exists \{G_j\}_{j=1}^\infty \subseteq  \mathrm{top}(\Rn) \quad \text{taki, że} \quad Y = \bigcap_{j=1}^\infty G_j.
              \end{equation*}
        \item Zbiór $Y$ nazywamy \textbf{zbiorem typu $F_\sigma$}, jeżeli jest on przeliczalną sumą zbiorów domkniętych, to znaczy:
              \begin{equation*}
                  \exists \, \{F_j\}_{j=1}^\infty \subseteq  \mathrm{cotop}(\Rn) \quad \text{taki, że} \quad Y = \bigcup_{j=1}^\infty F_j.
              \end{equation*}
    \end{enumerate}
\end{defi}
\begin{wn} W przestrzeni euklidesowej $\Rn$ (i ogólniej w dowolnej przestrzeni metrycznej):
    \begin{enumerate}
        \item Każdy zbiór domknięty jest typu $G_\delta$, tj. $\mathrm{cotop}(\Rn) \subseteq  G_\delta$.
        \item Każdy zbiór otwarty jest typu $F_\sigma$, tj. $\mathrm{top}(\Rn) \subseteq  F_\sigma$.
        \item Wszystkie zbiory klas $G_\delta$ oraz $F_\sigma$ są zbiorami borelowskimi ($G_\delta \subseteq  \BorelRn$ oraz $F_\sigma \subseteq  \BorelRn$).
    \end{enumerate}
\end{wn}
